% Template for post or presentation abstracts submitted to ILAS 2022
% 1. Fill in the presenter's information (name, affiliation, email
% address) and talk information (title of presentation, abstract).
% 2. Compile to make sure there are no errors.
% 3. Save the .tex file for your abstract with a distinctive name,
% ideally "FamilynameGivenname.tex".
% 4. Upload your .tex file to https://clr.ie/129557
%       
% * Please keep your LaTeX commands simple, and avoid the use of
%    user-defined macros.
% * Include details of co-authors and funding acknowledgements after the abstract.
% * Upload only the LaTeX file (with .tex extension).
% * Questions: Contact Niall Madden (Niall.Madden@NUIGalway.ie)

\documentclass[a4paper]{amsart} 
\usepackage{amssymb} 
\usepackage{hyperref}


\usepackage{amsfonts}
\usepackage{amsmath}
\usepackage[mathscr]{eucal}


\newcommand{\DD}{\mathfrak{D}}
\newcommand{\BB}{\mathfrak{B}}
\newcommand{\rk}{\mathrm{rank}}

\pagestyle{empty}
\date{} 

%% IGNORE THE NEXT 4 LINES
\makeatletter % 
\def\labelenumi{\theenumi.}
\def\theenumi{\@arabic\c@enumi}
\makeatother
%% STOP IGNORING NOW

\begin{document}

\title{The 2-Steiner distance matrix of a tree}
\author{Krishnan Sivasubramanian} %Presenting author only
\address{Department of Mathematics, IIT Bombay, Mumabi 400076, India} % Affiliation of presenting author
\email{krishnan@math.iitb.ac.in} % Email address of presenting author only

\maketitle

% The abstract  ***no more than one page***

Let $T$ be a tree with vertex set $V(T)=\{1,2,\ldots,n\}$. 
The Steiner distance of $S\subseteq V(T)$ of vertices
of $T$ is defined to be the number of edges in a smallest 
connected subtree of $T$ that contains all the vertices of $S$. 
The $k$-Steiner distance matrix $\DD_k(T)$ of $T$ is the 
$\binom{n}{k}\times \binom{n}{k}$ matrix whose rows 
and columns are indexed by subsets of vertices of size $k$.
The entry in the row indexed by $P$ and column indexed 
by $Q$ is equal to Steiner distance of $P \cup Q$. 
When $k=1$, it is easy to see that we have
$\DD_1(T) = D(T)$, the distance matrix of $T$ which is a 
well studied matrix.  In this work, we consider the next 
case, that is when $k=2$.  We show that $\rk(\DD_2(T)) = 2n -p - 1$ 
where $p$ is the number of pendant vertices (or leaves) in $T$.
We construct a basis $\BB$ for the row space of $\DD_2(T)$ and obtain 
a formula for the inverse of the nonsingular square submatrix 
$\DD=\DD_2(T)[\BB,\BB]$.  We also compute the determinant of 
$\DD$ and show that its absolute value is independent of the 
structure of $T$ and 
apply it to obtain the inertia of $\DD_2(T)$.  
\emph{This is joint work with Ali Azimi (Iran).}


%Please, follow the instructions at \url{http://ilas2022.ie/abstracts/}
%for submitting your abstract. Make sure you upload the \texttt{.tex}
%file, and not any others. 

%Try to keep your file as simple as possible, so that it is easily
%rendered in various formats. Give the file a name that readily
%identifies it with you, for example \texttt{OlgaTaussky.tex}. 
%
%If you wish to include citations, you can do so like this:
%\cite{my_key_for_OT49}. Use a bibitem key that is likely be be unique to
%you (e.g., don't use the default Math Reviews key).
%Any uncited bibliography entries will also
%appear. If you have none, then delete that section.

%% Coauthor details here (optional - delete if not applicable)
% and funding acknowledgment (optional - delete if not applicable)
%\emph{This is joint work with John Todd (Caltech) and Alexander Ostrowski (Basel).
%Supported by the National Mathematics Foundation, Grant XYZ-123.}

%\begin{thebibliography}{1}
%\bibitem{my_key_for_OT49}
%Olga Taussky.
%\newblock A recurring theorem on determinants.
%\newblock {\em  Amer. Math. Monthly}  56:672–676,  (1949).

%\bibitem{another_unique_key}
%Olga Taussky and John Todd.
%\newblock Cholesky, Toeplitz and the triangular factorization of
%symmetric matrices.
%\newblock {\em  Numer. Algorithms}, 41(2):197--202, 2006.
%\end{thebibliography}


\end{document}


% Please leave the next bit alone  
%%% Local Variables: 
%%% mode: latex
%%% TeX-master: "../ILAS-2022"
%%% End:
